\section{等距映射}
\label{sec:riem-isometries}

前面各节都在一个黎曼流形内部讨论。圆柱的例子已经提出了比较两个黎曼流形的问题：圆柱的一块与平面的一块“内在几何相同”，这句话要怎样准确表述？第~\ref{chap:manifold-relations}章用微分同胚表达两个流形具有相同的光滑结构，但微分同胚可以任意拉伸长度。这里要求映射在每个切空间上保持内积。

\paragraph{保持内积的条件}
设$(M,g)$与$(N,h)$是黎曼流形，$F:M\to N$光滑。映射的微分$dF_p:T_pM\to T_{F(p)}N$把切向量从$p$带到$F(p)$。我们要求：任取$p\in M$及$u,v\in T_pM$，
\begin{equation}
g_p(u,v)=h_{F(p)}\bigl(dF_pu,\,dF_pv\bigr).
\label{eq:riem-isometry-condition}
\end{equation}
右端是用$N$的度量量出的、被带过去的两个向量的内积，左端是原来的内积。在坐标中，若$F$的Jacobian矩阵为$\bm J$，条件就是$\bm G_M=\bm J^{\top}\bm G_N\bm J$；这与式\eqref{eq:riem-induced-metric}形式相同，诱导度量正是让包含映射满足此条件的度量。

\begin{definition}[黎曼等距映射与等距浸入]
\label{def:riem-isometry}
满足式\eqref{eq:riem-isometry-condition}的微分同胚$F:M\to N$称为\term{黎曼等距映射}（Riemannian isometry），此时称$(M,g)$与$(N,h)$等距。满足式\eqref{eq:riem-isometry-condition}的光滑映射称为\term{等距浸入}（isometric immersion）。若每点都有邻域使$F$限制在其上是到像集（$N$的开集）的等距映射，就称$F$为\term{局部等距}（local isometry）。
\end{definition}

条件\eqref{eq:riem-isometry-condition}本身就迫使$dF_p$单射：若$dF_pu=0$，则$g_p(u,u)=0$，从而$u=0$。所以满足该条件的光滑映射自动是浸入，定义中的名称与第~\ref{chap:manifold-relations}章一致。当$\dim M=\dim N$时，等距浸入的微分是可逆线性映射，由反函数定理它就是局部等距。

\paragraph{等距映射保持什么}
等距浸入保持每条曲线的长度。设$\gamma$是$M$中的分段光滑曲线，$F\circ\gamma$的速度是$dF(\dot\gamma)$，由式\eqref{eq:riem-isometry-condition}，$\lVert dF(\dot\gamma)\rVert_h=\lVert\dot\gamma\rVert_g$，所以$L_h(F\circ\gamma)=L_g(\gamma)$。对等距映射还可以反过来用$F^{-1}$，于是两边的连接曲线一一对应且长度相同，取下确界即得
\[
d_h\bigl(F(p),F(q)\bigr)=d_g(p,q).
\]
等距映射还把Levi-Civita联络带到Levi-Civita联络，因而保持平行移动、测地线和曲率。理由是唯一性：用$F$把$N$上的联络拉回到$M$上，得到的联络与$g$度量相容且无挠，所以就是$g$的Levi-Civita联络。这里省去逐条验证。

\paragraph{等距浸入可能缩短距离}
对等距浸入，只能得到一个方向的不等式。$M$中的曲线映成$N$中长度相同的曲线，但$N$中可能还有不来自$M$的更短路线，所以
\[
d_h\bigl(F(p),F(q)\bigr)\leq d_g(p,q).
\]
单位圆取弧长度量，包含映射$\mathbb S^1\hookrightarrow\mathbb R^2$是等距浸入：圆心角$\theta\leq\pi$的两点，内在距离为$\theta$，在平面中的距离是弦长$2\sin(\theta/2)$，严格更短。第~\ref{sec:riem-length-distance}节中“环境距离不超过内在距离”的结论，正是这个不等式的特例。

\begin{example}[把平面卷成圆柱]
\label{ex:riem-plane-cylinder}
映射$F:\mathbb R^2\to\{x^2+y^2=a^2\}\subseteq\mathbb R^3$，$F(s,z)=(a\cos(s/a),a\sin(s/a),z)$，由第~\ref{sec:riem-intrinsic-extrinsic}节的计算满足$\bm J^{\top}\bm J=\bm I$，所以平面取欧氏度量、圆柱取诱导度量时，$F$是局部等距。它不是单射：$(0,0)$与$(2\pi a,0)$映到同一点，平面中二者相距$2\pi a$，圆柱上距离为零。限制在带形$(0,2\pi a)\times\mathbb R$上，$F$是到圆柱去掉一条直线的等距映射。

这个例子说明了两点。局部等距保持每一小块的全部内在几何，包括曲率为零；但整体上它可以把远处的点粘到一起，距离不必保持。圆柱与平面局部等距，整体上不等距：圆柱上有绕一圈的不可收缩回路，平面上没有。
\end{example}

\begin{example}[球面的旋转]
\label{ex:riem-sphere-rotations}
正交矩阵$\bm Q$满足$\bm Q^{\top}\bm Q=\bm I$，它把单位球面映到自身，$F(x)=\bm Qx$，微分是$dF_x(v)=\bm Qv$。由$\langle\bm Qu,\bm Qv\rangle=\langle u,v\rangle$，条件\eqref{eq:riem-isometry-condition}成立，所以$F$是$\mathbb S^2$到自身的等距映射。它把大圆映成大圆，把例~\ref{ex:riem-sphere-transport}中的三角形回路映成另一个同样的回路，平行移动的转角不变。所有这样的映射组成一个群，即第~\ref{chap:lie-groups}章的正交群；把等距映射看成作用在流形上的变换，是那一章的话题。
\end{example}

\paragraph{与距离等距的关系}
从度量空间的角度，也可以只要求$d_h(F(p),F(q))=d_g(p,q)$对所有$p,q$成立，而不预先假设$F$光滑。对黎曼流形之间的满射，可以证明这样的映射自动光滑并满足条件\eqref{eq:riem-isometry-condition}，因而两种意义的等距一致；这一结论这里不需要，也不证明。本章使用的始终是由式\eqref{eq:riem-isometry-condition}定义的黎曼等距映射。
